% ============================================================
% Section 63: 简立方晶体的滑移面与临界切应力
% ============================================================
\section{固体理论：简立方晶体的滑移面与临界切应力}
\label{sec:slip-plane-critical-stress}

\begin{modelbox}[题目：滑移面判定与临界切应力的理论推导]
(1) 说明 $(0,0,1)$ 晶面是晶格常数为 $a$ 的简立方晶体中一组滑移面；

(2) 在外力作用下，简立方晶体 $(0,0,1)$ 晶面沿 $[1,0,0]$ 方向发生相对滑移，当滑移量为 $x$ 时，势能函数为
\begin{equation}
U = U_0 - A\cos\!\left(\frac{2\pi x}{a}\right),
\end{equation}
试用弹性理论证明临界切应力
\begin{equation}
\tau_m = \frac{G}{2\pi},
\end{equation}
其中 $G$ 为晶体的切变模量。
\end{modelbox}

\begin{tipbox}[考点快速分析]
\begin{itemize}
    \item \textbf{滑移面判据}：面间距越大，晶面间结合力越弱，越容易发生滑移
    \item \textbf{周期性势能}：晶体滑移的势能是位移 $x$ 的周期函数，周期为晶格常数 $a$
    \item \textbf{小应变极限}：$x\ll a$ 时，$\	au\propto x$，应与胡克定律 $\tau=G\gamma$ 一致
    \item \textbf{理论与实验差距}：理想晶体理论临界切应力 $\sim G/10$，实验值低 $3\sim 5$ 个数量级——这是\textbf{位错}存在的直接证据
\end{itemize}
\end{tipbox}

\begin{derivationbox}[标准解题过程]
\textbf{Step 1: $(0,0,1)$ 晶面作为滑移面的说明}

简立方晶体的正格子基矢为
\begin{equation}
\vec a_1 = a\,\hat{\imath},\qquad
\vec a_2 = a\,\hat{\jmath},\qquad
\vec a_3 = a\,\hat{k}.
\end{equation}
对应的倒格子基矢为
\begin{equation}
\vec b_1 = \frac{2\pi}{a}\,\hat{\imath},\qquad
\vec b_2 = \frac{2\pi}{a}\,\hat{\jmath},\qquad
\vec b_3 = \frac{2\pi}{a}\,\hat{k}.
\end{equation}

晶面族 $(hkl)$ 的晶面间距为
\begin{equation}
d_{hkl} = \frac{2\pi}{|\vec G_{hkl}|} = \frac{a}{\sqrt{h^2+k^2+l^2}}.
\end{equation}

对于 $(0,0,1)$ 面：
\begin{equation}
d_{001} = \frac{a}{\sqrt{0^2+0^2+1^2}} = a.
\end{equation}
对比其他晶面：$d_{110}=a/\sqrt{2}\approx 0.707a$，$d_{111}=a/\sqrt{3}\approx 0.577a$。$(001)$ 面具有\textbf{最大的面间距}。

\textit{结论}：晶面间距越大，相邻晶面间的原子结合力越弱，越容易发生相对滑移。因此 $(0,0,1)$ 面是简立方晶体的主要滑移面。

\textbf{Step 2: 临界切应力的理论推导}

\textit{势能与恢复力。}设 $(001)$ 晶面沿 $[100]$ 方向相对滑移量为 $x$。晶体的周期性决定了势能是 $x$ 的周期函数，周期为 $a$。取最简单的余弦近似：
\begin{equation}
U(x) = U_0 - A\cos\!\left(\frac{2\pi x}{a}\right).
\end{equation}
恢复力为
\begin{equation}
F = -\frac{\partial U}{\partial x} = -\frac{2\pi A}{a}\sin\!\left(\frac{2\pi x}{a}\right).
\end{equation}

设滑移面的面积为 $S$，则外加切应力 $\tau$ 应平衡恢复力产生的内应力：
\begin{equation}
\tau = -\frac{F}{S} = \frac{2\pi A}{aS}\sin\!\left(\frac{2\pi x}{a}\right)
= \tau_m\sin\!\left(\frac{2\pi x}{a}\right),
\end{equation}
其中最大切应力（即临界切应力）为
\begin{equation}
\tau_m = \frac{2\pi A}{aS}.
\label{eq:tau-m-A}
\end{equation}

\textit{小应变下的胡克定律。}当滑移量 $x\ll a$ 时，
\begin{equation}
\sin\!\left(\frac{2\pi x}{a}\right) \approx \frac{2\pi x}{a},
\end{equation}
切应力近似为
\begin{equation}
\tau \approx \tau_m\cdot\frac{2\pi x}{a}.
\end{equation}

切应变定义为滑移量 $x$ 与相邻晶面间距 $d$ 的比值：$\gamma = x/d$。对于 $(001)$ 面，$d=a$，故 $\gamma = x/a$。

根据胡克定律，切应力与切应变成正比：
\begin{equation}
\tau = G\gamma = G\,\frac{x}{a},
\end{equation}
其中 $G$ 为晶体的切变模量。

\textit{比较得出临界切应力。}将小应变近似式与胡克定律对比：
\begin{equation}
\tau_m\cdot\frac{2\pi x}{a} = G\,\frac{x}{a}.
\end{equation}
消去 $x/a$，得
\begin{equation}
\boxed{\tau_m = \frac{G}{2\pi} \approx 0.16\,G.}
\end{equation}
\textit{证毕。}
\end{derivationbox}

\begin{formulabox}[答案汇总]
\begin{center}
\begin{tabular}{ll}
\toprule
\textbf{问题} & \textbf{答案} \\
\midrule
(1) 滑移面判定 & $(0,0,1)$ 面是简立方中面间距最大的晶面（$d=a$），结合力最弱 \\
(2) 临界切应力 & $\displaystyle\tau_m = \frac{G}{2\pi} \approx 0.16\,G$ \\
\bottomrule
\end{tabular}
\end{center}
\end{formulabox}

\begin{insightbox}[物理图像：为什么理论强度远高于实验值？]
上述推导假设滑移是\textbf{晶面整体刚性平移}——同一时刻所有原子同时越过势垒。这要求外加应力足以克服整个晶面的结合力，故理论临界切应力高达 $G/2\pi\sim 0.16G$（更精细的估算给出 $\sim G/6$ 或 $\sim G/30$），即 $10^{10}\sim 10^{11}\,\mathrm{Pa}$ 量级。

然而实验测得的实际屈服强度远低于此：
\begin{itemize}
    \item 纯金属单晶：$\tau_{\mathrm{exp}}\sim 10^{-5}\,G$（如 Ag：$1/45000$）
    \item 即使是``软''金属（Pb、Al）：也只有 $10^{-4}\,G$
\end{itemize}

$3\sim 5$ 个数量级的差异说明：实际滑移绝非晶面整体平移，而是通过\textbf{位错}的逐步运动实现。

\textbf{位错机制}：位错线附近只有少量原子处于高位能畸变区，滑移时位错核心每次只移动几个原子间距，无需整个晶面同时翻越势垒。这就像移动一块大地毯：整体拖动极费力，但先推起一个褶皱再逐步推移则轻松得多。
\end{insightbox}

\begin{thinkbox}[考场速记：滑移问题的答题模板]
遇到``滑移面''或``临界切应力''类问题：
\begin{enumerate}
    \item \textbf{判滑移面}：计算各候选晶面的面间距 $d_{hkl}=a/\sqrt{h^2+k^2+l^2}$，取最大者
    \item \textbf{写周期势能}：$U(x)=U_0-A\cos(2\pi x/a)$（周期为晶格常数 $a$）
    \item \textbf{求恢复力与切应力}：$F=-\partial U/\partial x$，$\tau=-F/S=\tau_m\sin(2\pi x/a)$
    \item \textbf{取小应变极限}：$x\ll a$ 时线性化，$\tau\approx\tau_m(2\pi x/a)$
    \item \textbf{联立胡克定律}：$\tau=G\gamma=Gx/d$，注意 $d$ 是滑移面间距（$(001)$ 面 $d=a$）
    \item \textbf{比较得结果}：消去 $x$，得 $\tau_m=Gd/(2\pi a)$；若 $d=a$ 则 $\tau_m=G/(2\pi)$
    \item \textbf{讨论理论与实验差距}：提及位错理论，解释实际强度低 $3\sim 5$ 个数量级
\end{enumerate}
\textit{陷阱提醒}：切应变定义 $\gamma=x/d$ 中的 $d$ 是\textbf{相邻滑移面间距}，不是滑移方向的晶格常数。对于 $(001)$ 面沿 $[100]$ 滑移，$d=a$；但若滑移面是 $(111)$，则 $d=a/\sqrt{3}$，结果不同。
\end{thinkbox}
